\section{The Hybrid Monte Carlo Algorithm (HMC)}\label{sec: HMC}
The Hybrid Monte Carlo Algorithm (HMC) algorithm was first introduced in this paper: \cite{Duane1987HMC}. The authors solved the major
computational challenge of involving fermions into lattice field theory. \\\\
Simulating fermions in a lattice field theory is difficult, since fermions are described by Grassmann numbers, for which Monte Carlo methods
cannot be applied directly. Integrating over fermions produces very non-local effects in the action, which makes calculations extremely 
expensive. To solve this problem, the nested Monte Carlo steps are replaced with an equation of motion in a fictitious time variable. These are solved 
numerically with different techniques. Evolving the entire field configuration simultaneously using a numerical integrator avoids the locality problem 
but introduces truncation errors from discrete integration. Although these errors can be extrapolated away, they remain undesirable. 
The paper \cite{Duane1987HMC} introduces the HMC, which addresses the problem with the following steps: 
\begin{itemize}
    \item evolve all fields in parallel using a molecular-dynamics-like update, 
    \item a global Metropolis accept/reject step, which removes the truncation errors. 
\end{itemize}
This method combines deterministic molecular dynamics and accept/reject steps, achieving efficiency and exactness. Even though this algorithm 
was originally developed for lattice field theories, it is nowadays applied in various other fields. 
\\\\
The goal is again to calculate the expectation of some operator along with its partition function, as defined in the Equations 
\ref{eq: path_integral} and \ref{eq: partition_function}. The Metropolis algorithm is a rather inefficient way to do this and we search for 
a method to choose candidate configurations efficiently. We aim for a high acceptance rate and minimum correlation between successive configurations. 
For this, we introduce an artificial computer time $\tau$ and a hamiltonian dynamics to describe the development of the field $\phi(\tau)$. In addition, 
we will need the conjugate momenta $\pi(\tau)$. This means, instead of sampling the distribution $p(\phi) = \frac{1}{Z} e^{-S(\phi)}$, we sample 
$p(\phi, \pi) = \frac{1}{Z'} e^{-H[\phi, \pi]}$. The Hamiltonian $H[\phi, \pi]$ is given as follows: 
\begin{equation}
    H[\pi, \phi] = \frac{1}{2}\pi^2 + S(\phi), 
\end{equation}
where $S(\phi)$ is an arbitrary action. A key point of this algorithm is that the expectation values over $p(\phi)$ are the same as those over $p(\phi, \pi)$ 
up to some constant, since the integration over $\mathcal{D} \pi$ is a Gaussian integral, which gives rise to a constant. This constant will cancel out in the 
accept/reject step. \\\\
Introducing the conjugate momenta translates the numerical evaluation of the path integral into the language of classical Hamiltonian mechanics. The HMC algorithm 
solves the classical equations of motion for each lattice site $x$: 
\begin{align}\label{eq: HMC_DGLs}
    \dot{\phi}_x &= \frac{d \phi_x}{d \tau} = \frac{\partial H}{\partial \pi} = \pi_x \\
    \dot{\pi}_x &= \frac{d \pi_x}{d \tau} = -\frac{\partial H}{\partial \phi_x} = -\frac{\partial S}{\partial \phi_x} = -F[\phi_x] \qquad \text{"force"}
\end{align}
Because of the conservation of energy, the Hamiltonian remains constant. This can be used to check the accuracy of the solutions of the differential equations, 
which have to be solved numerically. In Chapter \ref{sec: integrators}, there will be a description of different methods to achieve this numerical solution.\\\\
The procedure for generating a new configuration $\phi'$ is to select some initial momenta $\pi$ from a standard Gaussian distribution (zero mean, unit variance)
\begin{equation}
    P_G(\pi) \propto \exp \left( -\frac{\pi^2}{2}\right), 
\end{equation}
and then evolve the system through the $(\phi, \pi)$- phase space for a time $\tau_0$ according to Hamilton's equations \ref{eq: HMC_DGLs}. 
The probability of choosing the configuration $(\phi', \pi')$ is 
\begin{equation}
    P_H((\phi, \pi) \rightarrow (\phi', \pi')) = \delta \left[ (\phi', \pi') - (\phi(\tau_0), \pi(\tau_0)) \right]. 
\end{equation}
The acceptance probability for this candidate is 
\begin{equation}
    P_A((\phi, \pi) \rightarrow (\phi', \pi')) = \min \left( 1, \exp( -\delta H) \right), 
\end{equation}
where $\delta H = H(\phi', \pi') - H(\phi, \pi)$. \\\\
In order for the HMC to correctly sample the target distribution $p(\phi) \propto e^{-S(\phi)}$, the transition probability restricted to the field variables $\phi$ alone must 
satisfy detailed balance. Starting from the full transition probability in $(\phi,\pi)$ space, the 
transition probability for $\phi \to \phi'$ is obtained by integrating out the momenta:
\begin{equation}
    P_M(\phi \to \phi') = \int [d\pi]\,[d\pi']\; P_G(\pi)\, P_H\!\big((\phi,\pi)\to(\phi',\pi')\big)\,P_A\!\big((\phi,\pi)\to(\phi',\pi')\big).
\end{equation}
Here $P_G(\pi)$ is the Gaussian distribution used to refresh the momenta, $P_H$ is the deterministic 
Hamiltonian evolution, and $P_A$ is the Metropolis accept/reject step.\\\\
The main requirement for detailed balance is that the molecular-dynamics evolution is reversible. 
This means that evolving $(\phi,\pi)$ forward in fictitious time $\tau$ to $(\phi',\pi')$ is equivalent to evolving $(\phi', -\pi')$ backward to $(\phi,-\pi)$:
\begin{equation}
    P_H\!\big((\phi,\pi)\to(\phi',\pi')\big) = P_H\!\big((\phi', -\pi')\to(\phi,-\pi)\big).
\end{equation}
This property holds automatically for Hamiltonian dynamics because the flow generated by a Hamiltonian is time-reversal invariant when accompanied by a sign flip of all momenta.
From the following identity (which we already used to show detailed balance for Metropolis)
\begin{align}
    e^{-H(\phi,\pi)}\min(1, e^{-\Delta H}) &= \min\!\left(e^{-H(\phi,\pi)}, e^{-H(\phi',\pi')} \right) \\
    &= e^{-H(\phi',\pi')} \min(1, e^{+\Delta H}),
\end{align}
the fact that the Hamiltonian $H(\phi,\pi)$ is also invariant under $\pi \to -\pi$ and $P_s(\phi) P_G(\pi) \propto \exp (-H(\phi, \pi))$, we obtain
\begin{equation}
    P_S(\phi)\, P_G(\pi)\,P_A((\phi,\pi)\to(\phi',\pi')) = P_S(\phi')\, P_G(\pi')\,P_A((\phi',\pi')\to(\phi,\pi)).
\end{equation}
Multiplying this by the transition probability $P_H$ and integrating over $\pi,\pi'$ gives
\begin{align}
    \int[d\pi]\,[d\pi']\;P_S(\phi)\,P_G(\pi)\,P_H((\phi,\pi)\to(\phi',\pi'))\, P_A((\phi,\pi)\to(\phi',\pi')) \\
    = \int[d\pi]\,[d\pi']\; P_S(\phi')\,P_G(\pi')\, P_H((\phi',\pi')\to(\phi,\pi))\, P_A((\phi',\pi')\to(\phi,\pi)).
\end{align}
Because the measure is invariant under momentum reversal,
\[
    [d\pi]\,[d\pi'] = [d(-\pi)]\,[d(-\pi')],
\]
and using the reversibility condition from above, this reduces precisely to
\begin{equation}
    P_S(\phi)\, P_M(\phi \to \phi') = P_S(\phi')\, P_M(\phi' \to \phi),
\end{equation}
which is the detailed balance condition for $P_M$.\\\\
By conservation of energy, we need $H = H'$, meaning $\delta H = 0$ and $P_A = 1$. However, solving differential equations numerically always introduces errors, hence
in general, $P_A \leq 1$. We have just proved that the HMC produces $\phi$ configurations with the correct distribution for any $\delta H \neq 0$. Also note that the momenta 
$\pi$ are not used to compute the observables. They can be disregared in each step \cite{Duane1987HMC}.\\\\
Note that the accept/reject step plays a crucial role in the HMC. Hamiltonian evolution exchanges energy between the kinetic term $\tfrac{1}{2}\pi^2$ and the potential term 
$S(\phi)$ while (ideally) keeping their sum $H = \tfrac{1}{2}\pi^2 + S(\phi)$ constant. The accept/reject step corrects the small energy violations $\delta H \neq 0$ resulting 
from the numerical integration, as it compares the Hamiltonians and accepts the proposed configuration with probability $P_A$. This enforces the correct equilibrium distribution, 
as a substantial increase in the energy will be rejected. In addition, the accept/reject step also ensures ergodicity. Pure Hamiltonian evolution is deterministic and reversible, 
and without momentum refreshment or a mechanism to leave a given energy surface, the algorithm would explore only a single trajectory in phase space. Periodically resampling the 
momenta from a Gaussian distribution redistributes kinetic energy, allowing the system to have different energies \cite{Duane1987HMC}. 